\section{Related Work}

We organise prior scholarship around three strands that together situate this study: (1) accessible video creation for blind and visually impaired (BVI) creators, where research prototypes have outpaced training infrastructure; (2) ability-diverse collaboration in HCI, where theoretical frames have shifted toward interdependence but empirical attention has centred on workplace settings rather than skill-learning; and (3) Universal Design for Learning and its critiques, which have been read separately from the empirical study of how access is enacted moment-to-moment.

\subsection{Accessible Video Creation: From Prototypes to Practice}

Personal media has become a core route through which BVI creators tell their own stories, build community identity, and challenge stereotypes \cite{emara_its_2025,seo_understanding_2021,brylla_communicative_2024}. Studies of BVI vloggers and bloggers find that creation supports wellbeing, expressive freedom, and visibility in digital spaces, even as platform accessibility remains uneven \cite{seo_understanding_2021,zhang_understanding_2023}. Despite these benefits, structured pathways into video creation are scarce; community surveys and creator studies report video as the most aspired-to and least supported skill area for BVI learners \cite{lyu_i_2024-1,zhang_understanding_2023}.

The barriers to BVI video creation are now well-documented across capture, editing, and publishing. A community study by Zhang et al. \cite{zhang_understanding_2023} found that nearly all BVI participants who had attempted video production encountered breakdowns at editing interfaces, visual feedback, and platform accessibility, with timeline-based interfaces, icon-driven controls, and preview windows particularly poorly supported by screen readers. Huh et al. \cite{huh_avscript_2023} synthesise four recurring challenges: inaccessible visual content, inability to verify visual quality, navigational opacity, and inconsistent screen-reader support. These constraints often push creators toward sighted assistance or restrict editing to the simplest operations \cite{lyu_i_2024-1,xiao_understanding_2025,raymond_examination_2024}.

Recent HCI prototypes propose alternative interaction paradigms that decouple editing from direct visual manipulation. AVscript \cite{huh_avscript_2023} replaces the timeline with a text-based interface that integrates transcripts, scene descriptions, and quality indicators, enabling navigation and selection by content rather than by frame. AltCanvas \cite{lee_altcanvas_2024} offers a tile-based grid for non-visual graphic design, providing accessible feedback for placement and composition. These systems demonstrate technical feasibility, but most remain at the proof-of-concept stage and are not publicly deployed. As a result, learners and practitioners must work with whatever publicly available tools admit partial accessibility, including iMovie, CapCut, and Videoleap, supplemented by AI-assisted aids such as Seeing AI for framing and visual description \cite{xiao_understanding_2025}. Self-taught creators have developed individual workarounds and shared knowledge through online communities \cite{lyu_i_2024-1,seo_understanding_2021}, yet these informal routes assume confidence and digital literacy that many beginners lack.

This split between demonstrated technical feasibility and limited deployment leaves a specific empirical gap. Less is known about how learners and ability-diverse teaching teams negotiate the friction between deployed mainstream tools and research prototypes when those tools sit at the centre of practical training. This study examines that question directly, taking the workshop as its analytical site rather than the individual interface.

\subsection{Ability-Diverse Collaboration in HCI}

A growing body of HCI scholarship reframes disability-related collaboration away from helper--helped dynamics. Wobbrock et al.'s \cite{wobbrock_ability-based_2011} ability-based design argues that systems should adapt to users' specific abilities rather than expecting users to conform to sighted templates, a position later extended to organisational and infrastructural settings \cite{nolte_implementing_2022}. Bennett et al.'s \cite{bennett_interdependence_2018} interdependence framework foregrounds the legitimate labour of coordination, mutual adjustment, and access work in mixed-ability collaboration, recasting access as collectively produced rather than individually granted. Mingus's \cite{mia_mingus_access_2017} concept of collective access offers a parallel formulation grounded in disability justice activism.

Empirical work has begun mapping the patterns through which ability-diverse collaboration operates. Xiao et al.'s \cite{xiao_systematic_2024} systematic review identifies recurrent collaboration patterns (ability sharing, ability combining) and technology roles (ability channel, ability supporter, ability combiner, communication supporter) across the literature. Studies of disabled professionals' collaborative ideation document persistent barriers in shared visual artefacts such as whiteboards and canvases, and argue for practices that externalise structure and surface access needs early \cite{das_that_2024}. Bandukda and Holloway \cite{bandukda_context_2024} examine how mixed-visual-ability dyads share information across context-dependent social interactions, showing that information needs and provision strategies shift moment-to-moment. Due et al.'s \cite{due_distributed_2021} concept of distributed perception describes how perception-related actions are shared across multiple semiotic agents in ability-diverse encounters, with embodied positioning playing a central role in enabling joint achievement.

Despite this progress, much of the empirical evidence draws from established workplace or social settings where roles are relatively stable and participants share prior context. Skill-learning settings present distinct conditions: novices, tutors, support workers, and co-learners occupy partially overlapping roles; the tools themselves are unfamiliar and partially accessible; and pedagogical sequencing constrains when and how access can be negotiated. Earlier studies of inclusive education have noted that group work risks reproducing dependence unless structure redistributes access labour \cite{ahmad_najmee_classroom_2025,barbieri_through_2025}, but few studies attend to the micro-interactional grain at which this redistribution succeeds or fails. This paper extends ability-diverse collaboration scholarship into the skill-learning workshop, asking what collaborative patterns emerge when the tool, the task, and the relationships are all simultaneously under construction.

\subsection{Universal Design for Learning and its Critiques}

Universal Design for Learning (UDL) is the dominant inclusive education paradigm in formal and community settings. Building on the principles of universal design in architecture, UDL specifies that learning environments should provide multiple means of representation, action and expression, and engagement, with the aim of accommodating learner variability without retrofitting \cite{rose_universal_2001,inc_udl_2024}. Empirical evaluations report gains in lesson planning, learner engagement, and instructional flexibility \cite{courey_improved_2013,rao_using_2016}. In BVI media education specifically, the UDL logic underpins arguments for multimodal scaffolds, alternative sensory channels, and accessible artefacts as standard provision rather than individual accommodation \cite{shaheen_exploring_2024,wilkens_using_2025,richard_jackson_audio-supported_2021}.

UDL has nonetheless drawn sustained critique from disability studies scholars who argue that universal design frameworks default to normative assumptions about users and contexts. Dolmage \cite{dolmage_academic_2017} traces how academic UDL reproduces ableist epistemologies even as it claims inclusion, by treating disability as an afterthought to be designed around rather than a generative starting point. Hamraie \cite{hamraie_building_2017} historicises universal design's roots in mid-twentieth-century white, middle-class user models and argues that ``access'' produced under such models embeds the very exclusions it claims to dissolve. These critiques converge on a structural point: UDL operates on a planning logic that presumes modalities can be specified in advance, and that the relevant user variability can be anticipated by designers. What it underspecifies is what happens when, in practice, the available modalities cannot meet a particular learner's needs in a particular moment.

Adjacent to UDL, a body of pedagogical work explores how inclusive maker and media education can be enacted with BVI learners. Hochberg et al.'s \cite{hochberg_supporting_2024} astronomy and computing curricula combine non-visual representations, shared inquiry, and teacher training for mixed-ability classes. Tactile and tangible approaches such as Khalaila and Lazar's \cite{khalaila_streetscape_2025} gamified navigation and Rocha et al.'s \cite{rocha_awareness_2025} tangible programming environments create shared workspaces for mixed-visual-ability collaboration, with explicit turn-taking and aligned metaphors. Sensory-substitution pedagogies extend this further, teaching ``visual'' concepts through structured audio and touch sequences \cite{he_charting_2024}. Across these studies, peer collaboration and co-teaching are repeatedly identified as essential supports for inclusive practice \cite{maesala_overcoming_2024,strogilos_meta-synthesis_2023,travers_collaboration_2024,miyauchi_systematic_2020}, yet the labour through which such supports are sustained, and the conditions under which they fail, remain under-examined.

Across these strands, UDL frames access as something that can be planned ahead of practice, while disability studies frames access as something collectively produced. What is missing is an empirical account of how access is actually enacted in moments where pre-arranged provision fails and learners, tutors, and helpers must construct workable interactional standards together. This paper contributes such an account through what we describe as dynamic standard-making, developed in the discussion in dialogue with the patterns we observe.
